\documentclass[aspectratio=169,11pt]{beamer}
\usepackage{teaching_slides}


\title{Multinomial Discrete Choice: Mixed Logit}
\author{Chris Conlon}
\institute{Grad IO}
\date{Fall 2026}

\begin{document}
\frame{\titlepage}


\begin{frame}{Mixed/ Random Coefficients Logit}
\small
An alternative is to allow for individuals to have \alert{random coefficients} $\beta_i$, mixing over various logits:
\begin{align*}
\beta_i &= \overline{\beta} + \Pi \, \textrm{y}_i  + \Sigma \, \nu_i
&\mu_{ij}=&   x_{j} \cdot \left(  \Pi \, \textrm{y}_i + \Sigma \, \nu_{i} \right) \\
u_{ijt} &= x_j \beta_i + \varepsilon_{ij}
&=&  x_j \beta  + \mu_{ij}+ \varepsilon_{ij}   \\
\sigma_{j}(\theta) &= \int \frac{\exp[x_{j} \beta_i ]}{1+\sum_k \exp[x_{k} \beta_i ]} f(\beta_i | \theta)
&=& \int \frac{\exp[x_{j} \beta + \mu_{ij} ]}{1+\sum_k \exp[x_{k} \beta+ \mu_{ik} ]} f(\symbf{\mu_i} | \theta)
\end{align*}
 \begin{itemize}
 \item $x_{j}$ are observed per usual, $\textrm{y}_i$ are individual demographics/features, and $\varepsilon_{ij}$ is IID Type I EV. The key is that $\nu_i$ is unobserved and mean zero (often normally distributed).
 \item Each individual draws a vector $\beta_i\symbf{/\mu_i}$ (separately from $\symbf{\varepsilon_i}$). Conditional on $\beta_i\symbf{/\mu_i}$ each person follows an IIA logit model.
 \item However we integrate (or mix) over many such individuals giving us a \alert{mixed logit} or \alert{heirarchical model} (if you are a statistician)
 \item It helps to think about fixing $\beta_i\symbf{/\mu_i}$ first and then integrating out over $\varepsilon_i$
 \end{itemize}
\end{frame}



\begin{frame}{Kinds of heterogeneity}
\begin{align*}
\sigma_{ij}(\theta) &= \int \frac{\exp[x_{j} \beta_i ]}{1+\sum_k \exp[x_{k} \beta_i ]} f(\symbf{\beta_i} | \theta) \approx \sum_{s=1}^S w_i^s \, \frac{\exp[x_{j} \beta_i^s ]}{1+\sum_k \exp[x_{k} \beta_i^s ]}
\end{align*}

\begin{itemize}
\item We can only \alert{approximate} the integral with \alert{weights} $w_{i}^s$ and \alert{nodes} $\beta_i^s$
\begin{itemize}
\item We can allow for there to be two types of $\beta_i$ in the population (high-type, low-type). \alert{latent class model}.
\item We can allow $\beta_i$ to follow an independent normal distribution for each component of $x_{ij}$ such as $\beta_i = \overline{\beta} + \Sigma \nu_i$.
\item We can allow for correlated normal draws $\beta_i \sim N(\mu_{\beta},\Sigma_{\beta})$.
\item Can allow for non-normal distributions too (lognormal, exponential).
\item Why is normal so easy?
\end{itemize}
\end{itemize}
\end{frame}

\begin{frame}{Adding Heterogeneity}
\begin{itemize}
\item The structure is extremely flexible but at a cost.
\item We generally must perform the integration numerically.
\item High-dimensional numerical integration is difficult. In fact, integration in dimension 8 or higher makes me very nervous.
\item We need to be parsimonious in how many variables have unobservable heterogeneity.
\item Again observed heterogeneity does not make life difficult so the more of that the better!
\begin{itemize}
\item If we see individual income, education, distance to hospital, etc. we can always interact that with observed characteristics without doing any additional integration.
\end{itemize}
\end{itemize}
\end{frame}


\begin{frame}
\frametitle{Mixed Logit}
How does it work?
 \begin{itemize}
\item Well we are mixing over individuals who conditional on $\beta_i$ or $\symbf{\mu_i}$ follow logit substitution patterns, however they may differ wildly in their $\sigma_{ij}$ and hence their substitution patterns.
\item For example if we are buying cameras: I may care a lot about price, you may care a lot about megapixels, and someone else may care mostly about zoom.
\item The basic idea is that we need to explain the heteroskedasticity of $Cov(\symbf{\varepsilon_i}, \symbf{\varepsilon_{\ell}},)$ what random coefficients do is let us use a basis from our $X$'s.
\item If our $X$'s are able to span the space effectively, then an RC logit model can approximate any arbitrary RUM (such as probit) (McFadden and Train 2002). 
\item Of course if you have 1000 products and two random coefficients, you are asking for a lot.
 \end{itemize}
\end{frame}

\begin{frame}
\frametitle{Elasticities}
At the level of an individual substitution patterns follow a plain logit:
\begin{align*}
\epsilon_{\sigma_j,p_j} = \frac{\partial \sigma_{j}}{\partial p_j} \cdot \frac{p_j}{\sigma_{j}}
=\frac{p_j}{\sigma_{j}}  \cdot \int \frac{\partial \sigma_{ij}}{\partial p_j}  \, d F_i 
= \frac{p_j}{\sigma_{j}}  \int \beta_i \cdot \sigma_{ij} \cdot (1-\sigma_{ij} ) \,d F_i 
\end{align*}
\begin{align*}
\epsilon_{\sigma_j,p_k} = \frac{\partial \sigma_{j}}{\partial p_k} \cdot \frac{p_k}{\sigma_{j}}
=\frac{p_k}{\sigma_{j}}  \cdot \int \frac{\partial \sigma_{ij}}{\partial p_k} \,  d F_i 
=- \frac{p_k}{\sigma_{j}}  \int \beta_i \cdot \sigma_{ij}\cdot \sigma_{ik}\, d F_i 
\end{align*}
Notation: $\sigma_{ij}(\symbf{\mu_i}) \equiv \sigma_{ij}$ is the individual choice probability and $\sigma_j = \int \sigma_{ij}\, d F_i$ is its aggregate (market-level) counterpart.
\end{frame}





\begin{frame}
\frametitle{Substitution (Conlon Mortimer 2020)}
At the level of an individual substitution patterns follow a plain logit:
\begin{align*}
D_{jk,i}(x) = \frac{\frac{\partial \sigma_{ik}}{\partial p_j}}{\left| \frac{\partial \sigma_{ij}}{\partial p_j} \right|} = \frac{\sigma_{ik}(x)}{1-\sigma_{ij}(x)}
\end{align*}
But at the aggregate level, we mix over heterogeneous individuals:
\begin{align*}
D_{jk}(x) = \int D_{jk,i}(x) w_i(z_j,z_j',x) d F_i \quad  \text{ where }
\quad w_i(z_j,z_j',x) = \frac{q_{ij}(z_j',x)-q_{ij}(z_j,x)}{q_{j}(z_j',x)-q_{j}(z_j,x)}
\end{align*}
Different interventions (price changes, product removals, quality changes, etc.) give different weighting schemes.
\end{frame}


\begin{frame}
\frametitle{Estimation Details: MSL/MSLE}
How do we estimate these models? (Derived from multinomial log-likelihood):
\begin{align*}
\theta_{MLE} = \arg \min_{\theta} - \sum_{i=1}^N \sum_{j=1}^J y_{ij} \ln \widehat{\sigma}_{ij}(\theta)
\end{align*}

\begin{itemize}
\item We need to perform numerical integration to get $\widehat{\sigma}_{ij}(\theta)$ and its derivative: $\widehat{\frac{\partial \sigma_{ij}}{\partial \theta}}$.
\item Consistency
\begin{itemize}
\item Technically for fixed number of MC draws the MSLE estimator is inconsistent. Why?
\item Very accurate integration is even more important!
\end{itemize}
\end{itemize}

\end{frame}

\begin{frame}
\frametitle{Estimation Alternatives: MSM}
Can also do a simulated GMM estimator with moment conditions:
\begin{align*}
\E\left[\left(y_{ij} - \widehat{\sigma}_{ij}(\theta) \right) | z_{ij} \right]=0
\end{align*}
\begin{itemize}
\item We need to perform numerical integration to get $\widehat{\sigma}_{ij}(\theta)$ and its derivative: $\widehat{\frac{\partial \sigma_{ij}}{\partial \theta}}$.
\item The \alert{optimal instruments} in the sense of Chamberlain (1987) or Amemiya (1977) are the derivative of moments with respect to parameters (\alert{scores}): $z_{ij} = \frac{\partial \log \sigma_{ij}}{\partial \theta}(\theta)$. (Much more on this when we get to BLP instruments next week.)
\item The true scores are infeasible (they depend on $\theta_0$).
\begin{itemize}
\item At true scores: same FOC as MLE.
\item Simulated scores have same consistency issue as MSLE.
\end{itemize}
\end{itemize}
\end{frame}


\begin{frame}
\frametitle{Estimation Alternatives: MSS}
As an alternative we could work with the score function directly:
\begin{align*}
\sum_{i=1}^n\widehat { \frac{\partial \log \sigma_{ij}}{\partial \theta}}(\theta) =\sum_{i=1}^N \frac{1}{\widehat{\sigma_{ij}}(\theta)} \cdot  \widehat{\frac{\partial \sigma_{ij}}{\partial \theta}}(\theta) = 0
\end{align*}
\begin{itemize}
\item We can get unbiased estimated of derivative (linearity)
\item We can't get unbiased estimate of $\frac{1}{\widehat{\sigma_{ij}}(\theta)}$
\begin{itemize}
\item Train's book suggests Accept-Reject 
\item Maybe importance sampling would work?
\end{itemize}
\end{itemize}
\end{frame}

\begin{frame}
\frametitle{Estimation Details: Hints}
 How bad is the simulation error?
 \begin{itemize}
\item Depends how small your shares are. 
\item Since you care about $\log \sigma_{jt}$ when shares are small, tiny errors can be enormous.
\item Often it is pretty bad.
%\item I recommend sticking with quadrature at a high level of precision.
%\item \texttt{sparse-grids.de} provide efficient high dimensional quadrature rules.
 \end{itemize}
\end{frame}

\begin{frame}
\frametitle{Estimation Details: Recommendations}
You should read these separate notes, but my recommendations are:
\begin{enumerate}
\item Numerical integration: 
\begin{itemize}
\item Use Gauss-Hermite quadrature rules to integrate over normal densities.
\item Quadrature is great in low dimensions -- but scales badly in high dimensions: if we need $N_a$ points to accurately approximate the integral in $d=1$ then we need $N_a^d$ points in dimension $d$ (using the tensor product of quadrature rules).
\item Follow Heiss and Winschel (JoE) and do \alert{sparse grids} in medium dimensions \item Try \url{http://sparse-grids.de}.
\end{itemize}
\item Nonlinear optimization: with mixtures the problem is non-convex (ie: very hard)
\begin{itemize}
\item You \alert{must} provide analytic gradients.
\item You should use a quasi-Newton solver and may need to try several.
\item Try multiple starting values.
\end{itemize}
\item Generalized Method of Moments
\begin{itemize}
\item You can read more about optimal IV in these settings.
\end{itemize}
\end{enumerate}
\end{frame}



\end{document}
